\documentclass{style/fulbright}
\usepackage{xeCJK}
\usepackage{bbm}
\usepackage{epigraph}
% Preamble fragment (utility math macros).
% =====================================================================
%  MATH MACROS  —  project-agnostic utility command set
%  File: <project>/style/math-macros.tex
% ---------------------------------------------------------------------
%  Usage in a paper's main.tex (after \documentclass and before
%  \begin{document}):
%
%      % =====================================================================
%  MATH MACROS  —  project-agnostic utility command set
%  File: <project>/style/math-macros.tex
% ---------------------------------------------------------------------
%  Usage in a paper's main.tex (after \documentclass and before
%  \begin{document}):
%
%      % =====================================================================
%  MATH MACROS  —  project-agnostic utility command set
%  File: <project>/style/math-macros.tex
% ---------------------------------------------------------------------
%  Usage in a paper's main.tex (after \documentclass and before
%  \begin{document}):
%
%      \input{style/math-macros}
%
%  Contains only domain-neutral commands that any econ/finance paper
%  in this vault could use. Project-specific symbol vocabulary
%  (\FTPL, \SDF, \HHwealth, ...) lives in notation-canon.tex.
%
%  Self-contained: assumes only what workingpaper.cls already loads
%  (amsmath, amsthm, bm, newtxmath). No new \usepackage calls.
% =====================================================================


% ---------------------------------------------------------------------
%  Probability and expectation operators
%
%  Convention: blackboard-bold E and P for measure-theoretic objects;
%  upright Var/Cov for moment statistics.
% ---------------------------------------------------------------------
\newcommand{\E}{\mathbb{E}}
\newcommand{\Prob}{\mathbb{P}}
\newcommand{\Var}{\mathrm{Var}}
\newcommand{\Cov}{\mathrm{Cov}}
\newcommand{\Corr}{\mathrm{Corr}}

% Conditional shortcuts: \Et[X], \Vart[X], \Covt[X,Y].
% (Subscript-as-token, so \Et^{\F_s} also works.)
\newcommand{\Et}{\E_{t}}
\newcommand{\Vart}{\Var_{t}}
\newcommand{\Covt}{\Cov_{t}}


% ---------------------------------------------------------------------
%  Number sets
% ---------------------------------------------------------------------
\newcommand{\R}{\mathbb{R}}
\newcommand{\N}{\mathbb{N}}
\newcommand{\Z}{\mathbb{Z}}
\newcommand{\Q}{\mathbb{Q}}
\newcommand{\Cset}{\mathbb{C}}        % \C is taken by some packages; alias
\newcommand{\Rplus}{\R_{+}}
\newcommand{\Rpp}{\R_{++}}            % strictly positive reals


% ---------------------------------------------------------------------
%  Distributions and convergence
% ---------------------------------------------------------------------
\newcommand{\Normal}{\mathcal{N}}                     % Gaussian
\newcommand{\Unif}{\mathcal{U}}                       % uniform
\newcommand{\iidsim}{\overset{\text{i.i.d.}}{\sim}}   % X_i \iidsim \Normal(0,1)
\newcommand{\todist}{\overset{d}{\longrightarrow}}    % convergence in distribution
\newcommand{\toprob}{\overset{p}{\longrightarrow}}    % convergence in probability
\newcommand{\toas}{\overset{a.s.}{\longrightarrow}}   % almost-sure convergence


% ---------------------------------------------------------------------
%  Differential calculus
%
%  \dd produces an upright "d" with thin-space spacing for integrals:
%      \int_0^1 f(x) \dd x
%
%  \pd / \pdtwo wrap the partial-derivative fraction:
%      \pd{f}{x}        ->  df/dx (partial)
%      \pdtwo{f}{x}     ->  d^2 f / dx^2  (partial)
%      \pdmix{f}{x}{y}  ->  d^2 f / dx dy  (mixed partial)
% ---------------------------------------------------------------------
\newcommand{\dd}{\mathop{}\!\mathrm{d}}
\newcommand{\pd}[2]{\frac{\partial #1}{\partial #2}}
\newcommand{\pdtwo}[2]{\frac{\partial^{2} #1}{\partial #2^{2}}}
\newcommand{\pdmix}[3]{\frac{\partial^{2} #1}{\partial #2 \, \partial #3}}


% ---------------------------------------------------------------------
%  Delimiters
%
%  Plain (non-auto-sized) versions. For auto-sizing in display math,
%  use \left\lvert ... \right\rvert directly, or load mathtools in
%  the paper's preamble and switch to \DeclarePairedDelimiter.
% ---------------------------------------------------------------------
\newcommand{\abs}[1]{\lvert #1 \rvert}
\newcommand{\norm}[1]{\lVert #1 \rVert}
\newcommand{\inner}[2]{\langle #1, #2 \rangle}
\newcommand{\set}[1]{\{ #1 \}}                        % set-builder shorthand
\newcommand{\given}{\;\middle|\;}                     % conditional bar in P(A \given B)


% ---------------------------------------------------------------------
%  Indicator function and common matrix ops
%
%  \argmax, \argmin, \plim are already declared in workingpaper.cls;
%  do not duplicate them here.
% ---------------------------------------------------------------------
\newcommand{\indic}{\mathbf{1}}                       % \indic\{X > 0\}
\DeclareMathOperator{\sign}{sign}
\DeclareMathOperator{\diag}{diag}
\DeclareMathOperator{\tr}{tr}
\DeclareMathOperator{\rk}{rank}
\DeclareMathOperator{\spn}{span}
\DeclareMathOperator{\Span}{span}                     % alias
\DeclareMathOperator{\proj}{proj}


% ---------------------------------------------------------------------
%  Vector and matrix typography (uses bm, loaded by workingpaper.cls)
%
%  Convention: \vect for column vectors, \mat for matrices. Both apply
%  bold-italic via \bm so that letters retain math-italic shape.
%  For upright bold (e.g. canonical R^n basis vectors), use \mathbf.
% ---------------------------------------------------------------------
\newcommand{\vect}[1]{\bm{#1}}
\newcommand{\mat}[1]{\bm{#1}}
\newcommand{\transpose}{^{\top}}                      % X\transpose -> X^T
\newcommand{\inv}{^{-1}}                              % X\inv      -> X^{-1}


% ---------------------------------------------------------------------
%  Spacing helpers
% ---------------------------------------------------------------------
\newcommand{\eqsep}{\;}                               % medium space inside equations
\newcommand{\where}{\quad\text{where }}               % "where" with spacing
\newcommand{\suchthat}{\;:\;}                         % shortcut for set-builder colon


\endinput

%
%  Contains only domain-neutral commands that any econ/finance paper
%  in this vault could use. Project-specific symbol vocabulary
%  (\FTPL, \SDF, \HHwealth, ...) lives in notation-canon.tex.
%
%  Self-contained: assumes only what workingpaper.cls already loads
%  (amsmath, amsthm, bm, newtxmath). No new \usepackage calls.
% =====================================================================


% ---------------------------------------------------------------------
%  Probability and expectation operators
%
%  Convention: blackboard-bold E and P for measure-theoretic objects;
%  upright Var/Cov for moment statistics.
% ---------------------------------------------------------------------
\newcommand{\E}{\mathbb{E}}
\newcommand{\Prob}{\mathbb{P}}
\newcommand{\Var}{\mathrm{Var}}
\newcommand{\Cov}{\mathrm{Cov}}
\newcommand{\Corr}{\mathrm{Corr}}

% Conditional shortcuts: \Et[X], \Vart[X], \Covt[X,Y].
% (Subscript-as-token, so \Et^{\F_s} also works.)
\newcommand{\Et}{\E_{t}}
\newcommand{\Vart}{\Var_{t}}
\newcommand{\Covt}{\Cov_{t}}


% ---------------------------------------------------------------------
%  Number sets
% ---------------------------------------------------------------------
\newcommand{\R}{\mathbb{R}}
\newcommand{\N}{\mathbb{N}}
\newcommand{\Z}{\mathbb{Z}}
\newcommand{\Q}{\mathbb{Q}}
\newcommand{\Cset}{\mathbb{C}}        % \C is taken by some packages; alias
\newcommand{\Rplus}{\R_{+}}
\newcommand{\Rpp}{\R_{++}}            % strictly positive reals


% ---------------------------------------------------------------------
%  Distributions and convergence
% ---------------------------------------------------------------------
\newcommand{\Normal}{\mathcal{N}}                     % Gaussian
\newcommand{\Unif}{\mathcal{U}}                       % uniform
\newcommand{\iidsim}{\overset{\text{i.i.d.}}{\sim}}   % X_i \iidsim \Normal(0,1)
\newcommand{\todist}{\overset{d}{\longrightarrow}}    % convergence in distribution
\newcommand{\toprob}{\overset{p}{\longrightarrow}}    % convergence in probability
\newcommand{\toas}{\overset{a.s.}{\longrightarrow}}   % almost-sure convergence


% ---------------------------------------------------------------------
%  Differential calculus
%
%  \dd produces an upright "d" with thin-space spacing for integrals:
%      \int_0^1 f(x) \dd x
%
%  \pd / \pdtwo wrap the partial-derivative fraction:
%      \pd{f}{x}        ->  df/dx (partial)
%      \pdtwo{f}{x}     ->  d^2 f / dx^2  (partial)
%      \pdmix{f}{x}{y}  ->  d^2 f / dx dy  (mixed partial)
% ---------------------------------------------------------------------
\newcommand{\dd}{\mathop{}\!\mathrm{d}}
\newcommand{\pd}[2]{\frac{\partial #1}{\partial #2}}
\newcommand{\pdtwo}[2]{\frac{\partial^{2} #1}{\partial #2^{2}}}
\newcommand{\pdmix}[3]{\frac{\partial^{2} #1}{\partial #2 \, \partial #3}}


% ---------------------------------------------------------------------
%  Delimiters
%
%  Plain (non-auto-sized) versions. For auto-sizing in display math,
%  use \left\lvert ... \right\rvert directly, or load mathtools in
%  the paper's preamble and switch to \DeclarePairedDelimiter.
% ---------------------------------------------------------------------
\newcommand{\abs}[1]{\lvert #1 \rvert}
\newcommand{\norm}[1]{\lVert #1 \rVert}
\newcommand{\inner}[2]{\langle #1, #2 \rangle}
\newcommand{\set}[1]{\{ #1 \}}                        % set-builder shorthand
\newcommand{\given}{\;\middle|\;}                     % conditional bar in P(A \given B)


% ---------------------------------------------------------------------
%  Indicator function and common matrix ops
%
%  \argmax, \argmin, \plim are already declared in workingpaper.cls;
%  do not duplicate them here.
% ---------------------------------------------------------------------
\newcommand{\indic}{\mathbf{1}}                       % \indic\{X > 0\}
\DeclareMathOperator{\sign}{sign}
\DeclareMathOperator{\diag}{diag}
\DeclareMathOperator{\tr}{tr}
\DeclareMathOperator{\rk}{rank}
\DeclareMathOperator{\spn}{span}
\DeclareMathOperator{\Span}{span}                     % alias
\DeclareMathOperator{\proj}{proj}


% ---------------------------------------------------------------------
%  Vector and matrix typography (uses bm, loaded by workingpaper.cls)
%
%  Convention: \vect for column vectors, \mat for matrices. Both apply
%  bold-italic via \bm so that letters retain math-italic shape.
%  For upright bold (e.g. canonical R^n basis vectors), use \mathbf.
% ---------------------------------------------------------------------
\newcommand{\vect}[1]{\bm{#1}}
\newcommand{\mat}[1]{\bm{#1}}
\newcommand{\transpose}{^{\top}}                      % X\transpose -> X^T
\newcommand{\inv}{^{-1}}                              % X\inv      -> X^{-1}


% ---------------------------------------------------------------------
%  Spacing helpers
% ---------------------------------------------------------------------
\newcommand{\eqsep}{\;}                               % medium space inside equations
\newcommand{\where}{\quad\text{where }}               % "where" with spacing
\newcommand{\suchthat}{\;:\;}                         % shortcut for set-builder colon


\endinput

%
%  Contains only domain-neutral commands that any econ/finance paper
%  in this vault could use. Project-specific symbol vocabulary
%  (\FTPL, \SDF, \HHwealth, ...) lives in notation-canon.tex.
%
%  Self-contained: assumes only what workingpaper.cls already loads
%  (amsmath, amsthm, bm, newtxmath). No new \usepackage calls.
% =====================================================================


% ---------------------------------------------------------------------
%  Probability and expectation operators
%
%  Convention: blackboard-bold E and P for measure-theoretic objects;
%  upright Var/Cov for moment statistics.
% ---------------------------------------------------------------------
\newcommand{\E}{\mathbb{E}}
\newcommand{\Prob}{\mathbb{P}}
\newcommand{\Var}{\mathrm{Var}}
\newcommand{\Cov}{\mathrm{Cov}}
\newcommand{\Corr}{\mathrm{Corr}}

% Conditional shortcuts: \Et[X], \Vart[X], \Covt[X,Y].
% (Subscript-as-token, so \Et^{\F_s} also works.)
\newcommand{\Et}{\E_{t}}
\newcommand{\Vart}{\Var_{t}}
\newcommand{\Covt}{\Cov_{t}}


% ---------------------------------------------------------------------
%  Number sets
% ---------------------------------------------------------------------
\newcommand{\R}{\mathbb{R}}
\newcommand{\N}{\mathbb{N}}
\newcommand{\Z}{\mathbb{Z}}
\newcommand{\Q}{\mathbb{Q}}
\newcommand{\Cset}{\mathbb{C}}        % \C is taken by some packages; alias
\newcommand{\Rplus}{\R_{+}}
\newcommand{\Rpp}{\R_{++}}            % strictly positive reals


% ---------------------------------------------------------------------
%  Distributions and convergence
% ---------------------------------------------------------------------
\newcommand{\Normal}{\mathcal{N}}                     % Gaussian
\newcommand{\Unif}{\mathcal{U}}                       % uniform
\newcommand{\iidsim}{\overset{\text{i.i.d.}}{\sim}}   % X_i \iidsim \Normal(0,1)
\newcommand{\todist}{\overset{d}{\longrightarrow}}    % convergence in distribution
\newcommand{\toprob}{\overset{p}{\longrightarrow}}    % convergence in probability
\newcommand{\toas}{\overset{a.s.}{\longrightarrow}}   % almost-sure convergence


% ---------------------------------------------------------------------
%  Differential calculus
%
%  \dd produces an upright "d" with thin-space spacing for integrals:
%      \int_0^1 f(x) \dd x
%
%  \pd / \pdtwo wrap the partial-derivative fraction:
%      \pd{f}{x}        ->  df/dx (partial)
%      \pdtwo{f}{x}     ->  d^2 f / dx^2  (partial)
%      \pdmix{f}{x}{y}  ->  d^2 f / dx dy  (mixed partial)
% ---------------------------------------------------------------------
\newcommand{\dd}{\mathop{}\!\mathrm{d}}
\newcommand{\pd}[2]{\frac{\partial #1}{\partial #2}}
\newcommand{\pdtwo}[2]{\frac{\partial^{2} #1}{\partial #2^{2}}}
\newcommand{\pdmix}[3]{\frac{\partial^{2} #1}{\partial #2 \, \partial #3}}


% ---------------------------------------------------------------------
%  Delimiters
%
%  Plain (non-auto-sized) versions. For auto-sizing in display math,
%  use \left\lvert ... \right\rvert directly, or load mathtools in
%  the paper's preamble and switch to \DeclarePairedDelimiter.
% ---------------------------------------------------------------------
\newcommand{\abs}[1]{\lvert #1 \rvert}
\newcommand{\norm}[1]{\lVert #1 \rVert}
\newcommand{\inner}[2]{\langle #1, #2 \rangle}
\newcommand{\set}[1]{\{ #1 \}}                        % set-builder shorthand
\newcommand{\given}{\;\middle|\;}                     % conditional bar in P(A \given B)


% ---------------------------------------------------------------------
%  Indicator function and common matrix ops
%
%  \argmax, \argmin, \plim are already declared in workingpaper.cls;
%  do not duplicate them here.
% ---------------------------------------------------------------------
\newcommand{\indic}{\mathbf{1}}                       % \indic\{X > 0\}
\DeclareMathOperator{\sign}{sign}
\DeclareMathOperator{\diag}{diag}
\DeclareMathOperator{\tr}{tr}
\DeclareMathOperator{\rk}{rank}
\DeclareMathOperator{\spn}{span}
\DeclareMathOperator{\Span}{span}                     % alias
\DeclareMathOperator{\proj}{proj}


% ---------------------------------------------------------------------
%  Vector and matrix typography (uses bm, loaded by workingpaper.cls)
%
%  Convention: \vect for column vectors, \mat for matrices. Both apply
%  bold-italic via \bm so that letters retain math-italic shape.
%  For upright bold (e.g. canonical R^n basis vectors), use \mathbf.
% ---------------------------------------------------------------------
\newcommand{\vect}[1]{\bm{#1}}
\newcommand{\mat}[1]{\bm{#1}}
\newcommand{\transpose}{^{\top}}                      % X\transpose -> X^T
\newcommand{\inv}{^{-1}}                              % X\inv      -> X^{-1}


% ---------------------------------------------------------------------
%  Spacing helpers
% ---------------------------------------------------------------------
\newcommand{\eqsep}{\;}                               % medium space inside equations
\newcommand{\where}{\quad\text{where }}               % "where" with spacing
\newcommand{\suchthat}{\;:\;}                         % shortcut for set-builder colon


\endinput


% Paper-local package additions (uncomment as needed).
% \usepackage{longtable}
% \usepackage{multirow}
% \usepackage{threeparttable}

% Paper-local masthead overrides (brand defaults live in the class).
\classification{Economics of Consumption}
% \subject{Course: ...}   % uncomment to add a subject line

\title{Research Proposal: Shaky Houses - A Dynamic Model of Housing Choice, House Prices, and Aggregate Fluctuations}
\author{Nguyen Chi Thang \\
  \small{Fulbright University Vietnam, Ho Chi Minh City, Vietnam \\
         \textit{thang.nguyen.240184@student.fulbright.edu.vn}}}
\date{\today}


\begin{document}
\maketitle

\epigraph{%
  \textit{各安其居而樂其業} \\
  \textit{các an kỳ cư nhi lạc kỳ nghiệp} \\
  Each is at ease in his dwelling and takes pleasure in his work.
  \\[1.2ex]
  \textit{安居樂業} \\
  \textit{an cư lạc nghiệp} \\
  To live in peace and work happily.%
}{Ban Gu (班固), \textit{Han Shu} (漢書)}

\section{Introduction}
The Chinese idiom “安居樂業” says when people live in peace, work happily, then there is a harmonious society. Such logic can be found in many cultures, where having a sheltered place to live is perceived as a precodition for a stable life, a stable society, and foremost, a stable economy. This makes housing services a special consumption of the household, obtainable via owning or renting, and more specifically, in case of owning, a household can decide whether it would finance the purchase by itself or a mortgage contract. Thus, I propose a simple and preliminary model to understand households' consumption behaviors, housing choices, and mortgage default in terms of aggregate shocks, with adopted frameworks from \citeA{campbell2015model,berger2018house,kaplan2020housing}. This model studies the how shocks materialized in three dynamics—(1) housing choices and financing options of the households, (2) housing market prices, and (3) aggregate measures. Main hypotheses are: \begin{enumerate}
  \item Randomly occurred aggregate shocks create changes in belief, therefore destablize housing market prices and aggregate measures of the economy.
  \item Housing markets prices and aggregate measures of the economy influence households' adjustments of belief, therefore changes in consumption.
  \item Changes in household consumption lead to more changes in housing market prices and aggregate measures of the economy.
\end{enumerate}
Hence, the assessment of the model and prospected hypotheses can contribute to the literature of macrofinance, where financial frictions (in this case, household's liquidity preferences,  mortgage constraints, etc.) amplifies aggregate shocks to the economy. From such studies, we can understand more about the role of housing services consumption and prices via the lens of household-level decision, which is a prominent topics in Asian countries such as Vietnam and China in term of land-led growth, Japan in term of housing market stagnation, as well as financial crisis could stem from it, as in the case of US sub-prime mortgage colapse that led to the Global Financial Crisis 2007-2008. 




\section{Methodology}
This is a model-based research, where I develop the theoretical model according to three literature. What they
share is the insight that housing is both consumption and collateral, so borrowing-constrained households determine how shocks transmit. The unique point is the link of each studies—\citeA{campbell2015model} about the financing and default decision of households, \citeA{berger2018house} about the pass-through from house prices to consumption, and \citeA{kaplan2020housing} about the equilibrium of prices and rents, with beliefs driving the boom and financial frictions amplifying the bust. The main contribution of this research is tying them together in a model to understand the dynamics between households, the housing market, and aggregates, carrying the core technicals of \citeA{campbell2015model}, adding on the aggregate and market mechanics of \citeA{berger2018house} according to the insights from \citeA{kaplan2020housing}. The formulation of the model will be minimal for ease of analysis, but also sufficient to be examined in three approaches: (1) finite-period representative household as the benchmark, (2) finite-period with cohorts, and (3) infinite-period with Aiyagari-type wealth-heterogeneous households.

For technical details, built on \citeA{campbell2015model} and extended with \citeA{diaz2008user}'s endogeneity of housing tenure choices, a household maximizes expected lifetime utility derived from nondurable consumption, housing services, and bequests to offspring, formally,
\begin{equation}
  \max_{\{c_t,\, f_t,\, h_t,\, \mathbbm{1}_{\text{rent},t}\}_{t=0}^{T}} \mathbb{E}_0 \left[ \sum_{t=0}^{T} \beta^t \left( \frac{c_t^{1-\gamma}}{1-\gamma}  + \theta \frac{(\mathbbm{1}_{\text{rent},t}f_t + (1-\mathbbm{1}_{\text{rent},t})h_t)^{1-\gamma}}{1-\gamma} \right)  + \beta^{T+1} b \frac{W_{T+1}^{1-\gamma}}{1-\gamma} \right]
\end{equation}
subject to its budget constraint, where $c_t$ is nondurable consumption, $f_t$ is housing services derived from rented housing stock, $h_t$ is housing services derived from owned housing stock, $\mathbbm{1}_{\text{rent},t}$ equals one if the household rents at date $t$ and zero if it owns, $W_{T+1}$ is terminal wealth, $\beta$ is the time discount factor, $\gamma$ is the coefficient of relative risk aversion, $\theta$ is the relative weight of housing services to nondurable consumption, and $b$ is the relative weight of utility from bequests to lifetime consumption. Housing services can be obtained by renting or owning, and owning can be financed with the household's own funds or with a mortgage subject to liquidity constraints. The choice among these options is governed by the relative taste for housing services, $\theta$, and the elasticity of substitution between housing and nondurable consumption, $1/\gamma$, which the household evaluates under its beliefs about future housing demand. Beliefs are formed from aggregate variables, including expected inflation, expected real interest rates, and expected income. A household also observes housing market prices, which are determined by the expected stream of future housing services valued in units of marginal utility, formally
\begin{equation}
  q_t = \mathbb{E}_t\left[\sum_{s=1}^{\infty}\beta^{s}\,\theta\,
\frac{\varrho f_{t+s}^{-\gamma} + (1-\varrho)h_{t+s}^{-\gamma}}{c_t^{-\gamma}} \right]
\end{equation}
where $q_t$ is the price of housing, $\varrho$ is the share of the housing stock that is rented, and $\beta$ is the discount factor. In the finite-period cases, the sum is truncated at the terminal date and the discounted terminal value of the housing stock is added. As liquidity-constrained households cannot always act on their own valuations, the price is pinned down by the unconstrained marginal household, and the transversality condition excludes a bubble component in the baseline model. Prices, aggregate measures and households' wealth and liquidity determine each household's rent-or-own choice, mortgage default decision, and consumption response, which closes the three dynamics described above.

The model might not be able to be solved analytically or numerically, so that simplification is needed to obtain adequate closed-form solutions. Then, similar to the approach of previous research, I aim to find validations of the model from US household-level datasets like the Panel Study of Income Dynamics and Consumer Expenditure Survey as well as supplementals (or imputations if needed) from the Survey of Consumer Finances and the Federal Reserve Economic Data. Following \citeA{berger2018house} and \citeA{kaplan2020housing}, the income process, mortgage rate, loan-to-value limits, and transaction costs are set externally, while the remaining parameters are calibrated to match data moments. The key parameters are (1) the relative taste for housing and the elasticity of substitution, matched to homeownership rates and housing expenditure shares, (2) the discount factor and bequest motive, matched to the wealth distribution and the marginal propensity to consume out of housing wealth, (3) mortgage frictions, namely the default cost and the foreclosure discount, matched to foreclosure rates and home equity, and (4) the perceived persistence of housing demand, matched to the price-to-rent ratio.

% =====================================================================
%  Bibliography
% =====================================================================
\bibliography{references}

\end{document}
